% !PRINTED 469
\weberpage{469}
%<<CONT

hence (10) gives:
\weq{12}{B\left(m\,\frac{A+C}{2} - (m+1)\,B + A\right) \equiv 0 \quad (\mathrm{mod}\ 8).}
If now $m \equiv 3\ (\mathrm{mod}\ 4)$, then $A + C \equiv 0\ (\mathrm{mod}\ 4)$ and the expression standing in the parenthesis in (12) is odd; hence (12) is satisfied \gesperrt{only} under the assumption that
\[
B \equiv 0 \quad (\mathrm{mod}\ 8).
\]
In order to avoid repetitions, we shall assume once and for all, in the theorems to be stated in what follows, that $\omega$ is the root of such an equation (3) of discriminant $-4m$ in which $A$ is odd and $B$ divisible by 8. With this the theorem is proved:

4. \gesperrt{If $m \equiv 3\ (\mathrm{mod}\ 4)$, then $f(\omega)^6$ is a class invariant.}

If further $m \equiv 5\ (\mathrm{mod}\ 8)$, the congruence (12) reduces to:
\[
B\left(\frac{5C - A}{2} + 2B\right) \equiv 0 \quad (\mathrm{mod}\ 8),
\]
and this condition is fulfilled if $B$ is divisible by 4, but it is not fulfilled if $B$ is divisible only by 2 and not by 4. From this it follows:

5. \gesperrt{If $m \equiv 5\ (\mathrm{mod}\ 8)$, then $f(\omega)^{12}$ is a class invariant.}

If $m \equiv 1\ (\mathrm{mod}\ 8)$, nothing can be concluded from the congruence (12). This agrees with the circumstance that in this case $\Phi_m(u, u)$ depends only on $u^8$. In order to reach similar results here as well, we apply the transformation of the second order. We take in (7):
\weq{13}{n \equiv 1 \quad (\mathrm{mod}\ 8)}
and set:
\[
u\,v = \pm \sqrt{2}^{3},
\]
i. e.:
\[
f\left(\frac{c + \partial\omega}{a}\right)^{3} f(\omega)^{3} = \pm \sqrt{2}^{3},
\]
or [§ 34, (18)]:
\weq{14}{f\left(\frac{c + \partial\omega}{a}\right)^{3} = \pm f\left(\frac{\omega - 1}{\omega + 1}\right)^{3}.}
The equation (7) thereby passes into an equation:
\weq{15}{\Phi_n\left(u, \frac{\pm\sqrt{2}^{3}}{u}\right) = 0,}
%CONT>>
